% Options for packages loaded elsewhere
\PassOptionsToPackage{unicode}{hyperref}
\PassOptionsToPackage{hyphens}{url}
\PassOptionsToPackage{dvipsnames,svgnames,x11names}{xcolor}
\documentclass[
  10pt,
  fontset=fandol]{ctexart}
\usepackage{xcolor}
\usepackage[margin=16mm]{geometry}
\usepackage{amsmath,amssymb}
\setcounter{secnumdepth}{-\maxdimen} % remove section numbering
\usepackage{iftex}
\ifPDFTeX
  \usepackage[T1]{fontenc}
  \usepackage[utf8]{inputenc}
  \usepackage{textcomp} % provide euro and other symbols
\else % if luatex or xetex
  \usepackage{unicode-math} % this also loads fontspec
  \defaultfontfeatures{Scale=MatchLowercase}
  \defaultfontfeatures[\rmfamily]{Ligatures=TeX,Scale=1}
\fi
\usepackage{lmodern}
\ifPDFTeX\else
  % xetex/luatex font selection
\fi
% Use upquote if available, for straight quotes in verbatim environments
\IfFileExists{upquote.sty}{\usepackage{upquote}}{}
\IfFileExists{microtype.sty}{% use microtype if available
  \usepackage[]{microtype}
  \UseMicrotypeSet[protrusion]{basicmath} % disable protrusion for tt fonts
}{}
\makeatletter
\@ifundefined{KOMAClassName}{% if non-KOMA class
  \IfFileExists{parskip.sty}{%
    \usepackage{parskip}
  }{% else
    \setlength{\parindent}{0pt}
    \setlength{\parskip}{6pt plus 2pt minus 1pt}}
}{% if KOMA class
  \KOMAoptions{parskip=half}}
\makeatother
\setlength{\emergencystretch}{3em} % prevent overfull lines
\providecommand{\tightlist}{%
  \setlength{\itemsep}{0pt}\setlength{\parskip}{0pt}}
\usepackage{amsmath,amssymb,mathtools,needspace}
\xeCJKDeclareCharClass{CJK}{"2032,"0400->"04FF,"2160->"216B,"2460->"2473,"25A0,"25C7,"25CB,"25CF}
\IfFontExistsTF{Arial Unicode MS}{\xeCJKsetup{AutoFallBack=true}\setCJKfallbackfamilyfont{\CJKrmdefault}{Arial Unicode MS}}{}
\setcounter{secnumdepth}{0}
\setlength{\emergencystretch}{2em}
\allowdisplaybreaks
\usepackage{bookmark}
\IfFileExists{xurl.sty}{\usepackage{xurl}}{} % add URL line breaks if available
\urlstyle{same}
\hypersetup{
  pdftitle={差分问题的可解性},
  colorlinks=true,
  linkcolor={Maroon},
  filecolor={Maroon},
  citecolor={Blue},
  urlcolor={Blue},
  pdfcreator={LaTeX via pandoc}}

\title{差分问题的可解性}
\author{}
\date{}

\begin{document}
\maketitle

\subsubsection{5.7.3 差分问题的可解性}\label{ux5deeux5206ux95eeux9898ux7684ux53efux89e3ux6027}

我们知道，通过自变量的适当变换可以消除线性方程（5.7.6）中的一阶导数项。下面仅就缺一阶导数项的方程来讨论这一问题。考察边值问题

\[
\begin{cases}
y''-q(x)y=r(x),\\
y(a)=\alpha,\quad y(b)=\beta.
\end{cases}
\tag{5.7.9}
\]

这里假定 \(q(x)\ge0\)，其对应的差分问题是

\[
\begin{cases}
\displaystyle\frac{y_{n+1}-2y_n+y_{n-1}}{h^2}-q_ny_n=r_n\quad(n=1,2,\cdots,N-1),\\[6pt]
y_0=\alpha,\quad y_N=\beta.
\end{cases}
\tag{5.7.10}
\]

\href{../../source-images/pdf-151.jpeg}{原页}

现在论证差分问题的可解性。由于式（5.7.10）是关于变量 \(y_n\)（\(n=0,1,2,\cdots,N\)）的线性方程组，要证明它的解存在唯一，只要证明对应的齐次方程组只有零解。为此，我们引进下述极值原理。

\textbf{定理 5.2} 对于一组不全相等的数 \(y_n\)（\(n=0,1,\cdots,N\)），记

\[
l(y_n)=\frac{y_{n+1}-2y_n+y_{n-1}}{h^2}-q_ny_n,\quad q_n\ge0\qquad(n=1,2,\cdots,N-1).
\]

假定 \(l(y_n)\ge0\)（\(n=1,2,\cdots,N-1\)），则 \(y_n\) 的正的最大值只能是 \(y_0\) 或 \(y_N\)；如果 \(l(y_n)<0\)（\(n=1,2,\cdots,N-1\)），则 \(y_n\) 的负的最小值只能是 \(y_0\) 或 \(y_N\)。

\textbf{证明} 用反证法。考察 \(l(y_n)\ge0\) 的情形，设 \(y_m\)（\(0<m<N\)）是正的最大值，即 \(y_m=\max_{0\le n\le N}y_n=M>0\)，且 \(y_{m-1}\) 和 \(y_{m+1}\) 中至少有一个小于 \(M\)，此时有

\[
l(y_m)=\frac{y_{m+1}-2y_m+y_{m-1}}{h^2}-q_my_m<\frac{M-2M+M}{h^2}-q_mM=-q_mM,
\]

由于 \(q_m\ge0,M>0\)，故由上式推出 \(l(y_m)<0\)，此与题设矛盾。

此外，\(l(y_n)<0\) 的情形可类似地进行讨论。证毕。

作为这一定理的推论有如下定理。

\textbf{定理 5.3} 差分问题式（5.7.10）的解存在并且是唯一的。

\textbf{证明} 只要证明对应的齐次方程组

\[
\begin{cases}
\displaystyle l(y_n)=\frac{y_{n+1}-2y_n+y_{n-1}}{h^2}-q_ny_n=0\quad(n=1,2,\cdots,N-1),\\[6pt]
y_0=y_N=0
\end{cases}
\]

只有零解。由于这里 \(l(y_n)=0\)（\(n=1,2,\cdots,N-1\)），由极值原理知，\(y_n\) 的正的最大值和负的最小值只能是 \(y_0\) 或 \(y_N\)，而按边界条件 \(y_0=y_N=0\)，故所有的 \(y_n\)（\(n=0,1,\cdots,N\)）全为零。证毕。

\begin{center}\rule{0.5\linewidth}{0.5pt}\end{center}

\subsection{相关原页校注（agent 补充）}\label{ux76f8ux5173ux539fux9875ux6821ux6ce8agent-ux8865ux5145}

以下校注来自本节覆盖的原页，可能也涉及同页相邻小节；原文照录与校正说明分开保留。

\subsubsection{原书 PDF 页 151 的校注}\label{ux539fux4e66-pdf-ux9875-151-ux7684ux6821ux6ce8}

\href{../pages/pdf-151.md}{完整原页转录} · \href{../../source-images/pdf-151.jpeg}{原页图像}

\begin{quote}
校注（agent 补充）：定理 5.2 的负最小值一侧，原书确实写作严格条件 \(l(y_n)<0\)；定理 5.3 却在 \(l(y_n)=0\) 时同时引用正最大值、负最小值结论。此处原文均照录。补足该引用的方法是对 \(-y_n\) 应用定理 5.2 的 \(l\ge0\) 一侧（利用 \(l\) 的线性），从而排除负最小值；无需改动原书方程。
\end{quote}

\end{document}
